\section{Fixed-target extremizers and contact}
\label{v9:core:section}

We optimize the excess above a fixed candidate constant. This gives a
variational problem whose coefficient is known before any acquisition
step. The positive excess is retained through the contact calculation
and the conserved identity; it becomes a nonnegative term in the final
resource. No upper bound for the maximizing normalized score is needed.
Throughout, sums have
independent identically distributed summands with mean zero, variance one
and finite third absolute moment $\rho$. Fix $n\ge2$, $c>0$, and a real
raw sum threshold $H$. Set
\begin{gather*}
 z=H/\sqrt n,\quad F=\Pp(S_n\le H),\quad
 a=\sqrt n\{F-\Phi(z)\}-c\rho,\\
 \varphi=\phi(z),\quad\beta=\E[X|X|],\quad\gamma=\E X^3.
\end{gather*}
In particular $a$ is not divided by $\rho$.

\subsection{Attainment, contamination and the entire closed support}

\begin{proposition}\label{v9:core:attainment}
For each fixed $n$ and $c>0$, the supremum of $a$ over these laws and
thresholds is attained, and the closed-CDF discrepancy at every attaining
pair is positive.
At every maximizing pair,
\begin{align}
 I(y)={}&n^{3/2}\{h(y)-F\}+ny\varphi
       +\tfrac12H\varphi(y^2-1)\notag\\
 &+\tfrac32c\rho(y^2-1)-c\{|y|^3-\rho-3y\beta\}\le0,
 \qquad h(y)=\Pp(S_{n-1}\le H-y),\label{v9:core:contact}
\end{align}
with equality on the whole closed support. That support is compact.
Neither statement assumes $a\ge0$ or $H\ge0$.
\end{proposition}
\begin{proof}
Write $A$ for the supremum. Since $\rho\ge1$ and the discrepancy is
at most one, $A\le\sqrt n-c$. Symmetric Bernoulli evaluated at $H=n$
gives $A\ge\sqrt n(1-\Phi(\sqrt n))-c>-c$.
For a maximizing sequence, $a_j\to A$ and
$c\rho_j\le\sqrt n-a_j$ give a common third-moment ceiling.
The discrepancies satisfy $\sqrt n(F_j-\Phi(H_j/\sqrt n))
=a_j+c\rho_j\ge a_j+c>\eta$ eventually for some $\eta>0$.
Chebyshev bounds the negative threshold tail by $n/H_j^2$;
the positive discrepancy is at most the normal tail, also at most
$n/H_j^2$. Thus $H_j$ is bounded.

Tightness gives a subsequence $P_j\Rightarrow P$, $H_j\to H$ and
$\rho_j\to b$. The common third-moment ceiling $M$ gives tail
bounds $\int_{|x|>K}x^2dP_j\le M/K$ and
$\int_{|x|>K}|x|dP_j\le M/K^2$. Bounded continuous truncations
therefore preserve mean and variance. Lower semicontinuity gives
$1\le\rho(P)\le b$. For fixed $n$, the sum laws converge weakly.
Applying the closed-set inequality first at $H+t$ and then letting
$t\downarrow0$ gives
$\limsup F_j(H_j)\le F_P(H)$, including an atom at $H$.
Consequently
\[
 a(P,H)\ge A+c\{b-\rho(P)\}\ge A.
\]
Maximality forces equality and $\rho(P)=b$. Convergence of third
moments is a conclusion here; it was not inferred from bounded third
moments. The discrepancy remains positive since $A+c\rho\ge A+c>0$.

The contact inequality is a first-variation condition. To obtain it,
contaminate the raw marginal by
$(1-\epsilon)P+\epsilon\delta_y$. Its mean is $\epsilon y$,
variance $v_\epsilon=1+\epsilon(y^2-1)-\epsilon^2y^2$, and centered
third absolute moment has derivative $|y|^3-\rho-3y\beta$ at zero.
The last assertion follows by dominated differentiation of
$|x-\epsilon y|^3$, with bound $3|y|(|x|+|y|)^2$.
Evaluate the standardized law at the standardized image of the fixed
raw event $S_n\le H$. The raw event probability is a polynomial in
$\epsilon$, with derivative $n(h(y)-F)$ even at an atom collision.
The normal argument is
$(H-n\epsilon y)/\sqrt{nv_\epsilon}$, whose derivative is
$-\sqrt n y-H(y^2-1)/(2\sqrt n)$.
The one-sided derivative of the excess is exactly \eqref{v9:core:contact};
maximality makes it nonpositive.

Its integral against $P$ is zero. All terms are integrable.
The closed-event map $h$ is upper semicontinuous; at an atom it takes
the larger one-sided value. Thus $I$ is upper semicontinuous. Strict
negativity at a closed-support point would hold with a uniform negative
margin on a neighborhood of positive $P$-mass, contradicting its zero
integral. Finally $I(y)\le-c|y|^3+C_2y^2+C_1|y|+C_0$ for constants
depending on this pair. It is strictly negative on both sufficiently
distant rays, proving compact support. No higher moments were used
before this step.
\end{proof}

\subsection{Reflection, equality, and the first unresolved index}
\label{v9:core:reflection}
The sign of the attained maximum detects failure of a strict row bound,
including equality. If $C_n$ denotes the supremum of the normalized
absolute discrepancy, reflection gives, including an atom at $H$,
\[
 \Pp(-S_n\le-H)-\Phi(-z)
 =\Phi(z)-\Pp(S_n<H)\ge\Phi(z)-F.
\]
Thus each negative closed-CDF error is bounded by a positive closed-CDF
error of the reflected law; equality at an atom is not assumed.
Consequently $C_n$ is also the supremum of positive
closed scores $r$. Since $a=\rho(r-c)$, $C_n<c$ gives
$A\le C_n-c<0$, and $C_n>c$ gives $A>0$. If $C_n=c$, choose
$r_j\uparrow c$ with $r_j\ge c/2$. Then
$\rho_j\le2\sqrt n/c$, so $\rho_j(r_j-c)\to0$ and $A=0$.
Hence $C_n\ge c$ supplies an attained pair with $a\ge0$.

\begin{corollary}[First failure of a strict bound]\label{v9:core:first-failure}
Assume a genuine finite base $C_j\le c$ for $1\le j<N$. If a strict
bound $C_n<c$ fails at any index $n\ge N$, there is a least such index,
an attained fixed-target pair with $a\ge0$, and true predecessor bounds
$C_j\le c$ for every $j<n$. The bounds are strict for $N\le j<n$.
\end{corollary}
\begin{proof}
The nonempty set of failing indices at least $N$ has a least member.
The preceding equivalence supplies its maximizing pair. Indices below
$N$ use the finite base; the other earlier indices use minimality.
\end{proof}
Thus an argument excluding every such pair for $n\ge N$ would give
strict bounds above the finite base, including exclusion of equality.
The subsequent entry argument uses only $C_{n-1}\le c$; it does not
use an upper bound for the actual score $c+a/\rho$. The finite base
and the eventual exclusion remain separate obligations.

Equivalently, $\Delta=Q-h\ge0$ everywhere and $\Delta=0$ on the closed
support, where
\begin{equation}\label{v9:core:Q}
 Q(y)=F+\frac{H\varphi+c\rho}{2n^{3/2}}
 -\frac{n\varphi+3c\beta}{n^{3/2}}y
 -\frac{H\varphi+3c\rho}{2n^{3/2}}y^2
 +\frac c{n^{3/2}}|y|^3.
\end{equation}
